% Minimal retained preamble for lane-p2 cards taken from
% sources/logic-fl05-ramsey-like-theorems-for-the-schre/pinned-original-source.tex
% (arXiv:2412.11598v3, CC BY 4.0).  This source ships no separate preamble file:
% the pinned tex is a concatenation of the authors' main.tex and macros.tex and
% carries its preamble inline, so the blocks below are sliced mechanically from
% that file and nothing has been renamed or redefined.
%
% Copied verbatim from the source:
%   line 2        : the authors' class line
%   lines 3-5,7,9 : graphicx, amsthm, the ams bundle, cleveref, url
%   line 10       : tikz-cd, used by the authors' commutative diagrams
%   lines 20-108  : the authors' macro definitions, taken in full so that every
%                   macro of the retained spans is available
%   lines 113-124 : the authors' theorem declarations
%   lines 126-128 : the authors' macros for the three calligraphic letters, which
%                   sit just after the theorem declarations in the source
%   lines 131-141 : the authors' machinery for restating a theorem
%   lines 143-148 : the authors' coloured comment macros
%
% Omitted source preamble lines and why:
%   line 1, 110-111 and 150 (concatenation comment, title, date, author): the
%       assembler supplies the title block, author list and date;
%   line 6 (hyperref) and line 8 (inputenc with the utf8x option): the forced
%       xelatex engine supplies the hyperref wiring itself and utf8x is
%       incompatible with it;
%   line 115 (numberwithin theorem by section): the assembler numbers the
%       retained material itself;
%   lines 120 and 121-123 are declarations that the authors left commented out or
%       that duplicate earlier ones and are not reproduced.
\documentclass{article}
\usepackage{graphicx}
\usepackage{amsthm}
\usepackage{amsfonts, amsmath, amssymb, amsthm}
\usepackage{cleveref}
\usepackage{url}
\usepackage{tikz-cd}
\newcommand{\Psf}{\mathsf{P}}
\newcommand{\Qsf}{\mathsf{Q}}
\newcommand{\Rsf}{\mathsf{R}}
\def\F{\mathcal{F}}
\def\U{\mathcal{U}}
\def\C{\mathcal{C}}
\def\I{\mathcal{I}}
\def\D{\mathcal{D}}
\def\L{\mathcal{L}}
\def\V{\mathcal{V}}
\def\A{\mathcal{A}}
\def\B{\mathcal{B}}
\def\P{\mathcal{P}}
\def\M{\mathcal{M}}
\def\Nc{\mathcal{N}}
\def\Q{\mathcal{Q}}
\newcommand{\cs}{2^{\NN}}
\newcommand{\bs}{\NN^{\NN}}
\newcommand{\bstr}{2^{<\NN}}
\newcommand{\uh}{\upharpoonright}
\newcommand{\converge}{\!\!\downarrow}
\newcommand{\diverge}{\!\!\uparrow}
\newcommand{\qvdash}{\operatorname{{?}{\vdash}}}
\newcommand{\nqvdash}{\operatorname{{?}{\nvdash}}}
\newcommand{\ISig}{\mathsf{I}\Sigma^0}
\newcommand{\Robinson}{\mathsf{Q}}
\newcommand{\IPi}{\mathsf{I}\Pi^0}
\newcommand{\IDelta}{\mathsf{I}\Delta^0}
\newcommand{\Iopen}{\mathsf{I_{open}}}
\newcommand{\BSig}{\mathsf{B}\Sigma^0}
\newcommand{\BPi}{\mathsf{B}\Pi^0}
\newcommand{\BDelta}{\mathsf{B}\Delta^0}
\newcommand{\BCSig}{\mathsf{BC}\Sigma^0}
\newcommand{\BCPi}{\mathsf{BC}\Pi^0}
\newcommand{\BCDelta}{\mathsf{BC}\Delta^0}
\newcommand{\LSig}{\mathsf{L}\Sigma^0}
\newcommand{\LPi}{\mathsf{L}\Pi^0}
\newcommand{\CSig}{\mathsf{C}\Sigma}
\newcommand{\RCA}[0]{\mathsf{RCA}}
\newcommand{\WKL}[0]{\mathsf{WKL}}
\newcommand{\WWKL}[0]{\mathsf{WWKL}}
\newcommand{\ACA}[0]{\mathsf{ACA}}
\newcommand{\ATR}[0]{\mathsf{ATR}}
\newcommand{\PIOOCA}[0]{\Pi^1_1\mathsf{-CA}}
\newcommand{\CAC}[0]{\mathsf{CAC}}
\newcommand{\ADS}[0]{\mathsf{ADS}}
\newcommand{\RRT}[0]{\mathsf{RRT}}
\newcommand{\FS}[0]{\mathsf{FS}}
\newcommand{\PFS}[0]{\mathsf{PFS}}
\newcommand{\TS}[0]{\mathsf{TS}}
\newcommand{\GRRT}[0]{\mathsf{GRRT}}
\newcommand{\LRRT}[0]{\mathsf{LRRT}}
\newcommand{\EM}[0]{\mathsf{EM}}
\newcommand{\SEM}[0]{\mathsf{SEM}}
\newcommand{\HalfEM}[0]{\mathsf{Half}\mbox{-}\mathsf{EM}}
\newcommand{\RT}[0]{\mathsf{RT}}
\newcommand{\SRT}[0]{\mathsf{SRT}}
\newcommand{\NN}[0]{\mathbb{N}}
\newcommand{\nat}[0]{\mathbb{N}}
\newcommand{\card}{\operatorname{card}}
\newcommand{\finsub}{\subseteq_{\mathtt{fin}}}
\newcommand{\W}{\mathrm{W}}
\newcommand{\sW}{\mathrm{sW}}
\newcommand{\lex}{\mathrm{lex}}
\newcommand{\base}{\operatorname{base}}
\newcommand{\range}{\operatorname{ran}}
\def\qt#1{``#1''}%

\newtheorem*{statement}{Statement}
\newtheorem{theorem}{Theorem}
\newtheorem{maintheorem}[theorem]{Main Theorem}
\newtheorem{lemma}[theorem]{Lemma}
\newtheorem{question}[theorem]{Question}
\newtheorem{proposition}[theorem]{Proposition}
%\newtheorem{statement}[theorem]{Statement}
\newtheorem{remark}[theorem]{Remark}
\newtheorem{definition}[theorem]{Definition}
\newtheorem{corollary}[theorem]{Corollary}
\newtheorem{example}[theorem]{Example}
\newcommand{\R}{\mathcal{R}}
\renewcommand{\S}{\mathcal{S}}
\newcommand{\T}{\mathcal{T}}

% Allow to repeat theorems
\makeatletter
\newtheorem*{rep@theorem}{\rep@title}
\newcommand{\newreptheorem}[2]{%
\newenvironment{rep#1}[1]{%
 \def\rep@title{#2 \ref{##1}}%
 \begin{rep@theorem}}%
 {\end{rep@theorem}}}
\makeatother

\newreptheorem{theorem}{Theorem}

% Comments
\usepackage{xcolor}
\newcommand{\ludovic}[1]{\textcolor{blue}{#1}}
\newcommand{\quentin}[1]{\textcolor{green}{#1}}
\newcommand{\lorenzo}[1]{\textcolor{orange}{#1}}
\newcommand{\oriola}[1]{\textcolor{red}{#1}}
\crefname{originallemma}{original lemma}{original lemmas}
\Crefname{originallemma}{Original lemma}{Original lemmas}
\makeatletter
\newlabel{lem:fs-forcing-question-spec}{{6.24}{source}{}{sourcepointer-1}{}}
\newlabel{lem:fs-forcing-question-spec@cref}{{[originallemma][0][]6.24}{[1][0][]source}{}{}{}}
\makeatother
\crefname{originallemma}{original lemma}{original lemmas}
\Crefname{originallemma}{Original lemma}{Original lemmas}
\makeatletter
\newlabel{lem:fs-forcing-question-sigma01}{{6.23}{source}{}{sourcepointer-2}{}}
\newlabel{lem:fs-forcing-question-sigma01@cref}{{[originallemma][0][]6.23}{[1][0][]source}{}{}{}}
\makeatother
\title{The diagonalisation lemma: a condition can be extended so that the generic set differs from the fixed set everywhere}
\author{Lorenzo Carlucci, Oriola Gjetaj, Quentin Le Houérou, Ludovic Levy Patey}
\date{arxiv:2412.11598v3; CC BY 4.0}
\begin{document}
\maketitle
\noindent\textit{Editorial note.} The counted unit is the article's own Lemma lem:fs-diagonalization at source lines 1240-1243 together with its proof at source lines 1244-1251, a straight adjacent pairing in the pinned file. The pinned file is a concatenation of the authors' main and macro files; it is the article's fixed version arXiv:2412.11598v3 under CC BY 4.0.\par\medskip
\noindent\textit{Editorial note.} The retained preamble is the lane's shared slice for this source, 137 lines, carrying the authors' macro block together with the three calligraphic letter macros that sit between the macro block and the theorem declarations.\par\medskip
\noindent\textit{Editorial note.} Two context spans carry the notions the retained statement and proof work with: the definition of the forcing question at lines 1196-1199, retained from its opening marker through its closing marker, and the definition of forcing for a condition at line 1218. They are context, not counted units, and neither overlaps the counted interval.\par\medskip
\noindent\textit{Editorial note.} The retained proof names two lemmas of the same section, at source lines 1201-1204 and 1221-1227, both counted by separate cards of this lane and both declared here as non-load-bearing navigation pointers with the article's own section-wise numbers 6.23 and 6.24. The lemma on extensibility is named only in the surrounding prose just after the counted span and is therefore not a pointer of this card.\par\medskip
\noindent\textit{Editorial note.} The third case of the proof closes by a contradiction with the defining property of a condition; that property is quoted in the card's standard prerequisites and the context spans show the reservoir condition it refers to. No citation occurs in the retained spans, so no bibliography entry is supplied.\par\medskip
\noindent\textit{Allowed standard prerequisites (editorial).}\begin{itemize}
\item Conditions, the forcing question, forcing and the two effectiveness lemmas. A condition is a triple of a colouring constrained by a computable bounding function, a finite stem and an infinite reservoir computing a fixed set D; the forcing question holds for a condition and a formula when every relation of the bounded space admits a finite free block in the reservoir making the formula true on the enlarged stem, and the section proves that this question is first-order over one level of bounded progressive functions and that it meets its specification, that is, it tells apart the extensions that force the formula from those that force its negation. The author's own definitions and those two lemmas are the allowed inputs.. The counted lemma is the diagonalisation step: it is what makes the generic set differ from the fixed set, and the section's final theorem on strong cone avoidance invokes it immediately after the extensibility lemma.
\end{itemize}\medskip
\ifdefined\hypertarget\hypertarget{sourcepointer-1}{}\fi
\noindent\textit{Source reference.} Original lemma 6.24: sources/logic-fl05-ramsey-like-theorems-for-the-schre/pinned-original-source.tex, source lines 1221--1227. Navigation pointer invoked twice by the retained proof, once in each of the first two cases, to pass from a statement about the forcing question to an extension that forces the functional's value or its negation. That lemma is counted by a separate card of this lane, so its proof is not reproduced here. The number is the article's own section-wise number for it, counted over the numbered environments of section 6.\par\medskip
\ifdefined\hypertarget\hypertarget{sourcepointer-2}{}\fi
\noindent\textit{Source reference.} Original lemma 6.23: sources/logic-fl05-ramsey-like-theorems-for-the-schre/pinned-original-source.tex, source lines 1201--1204. Navigation pointer invoked by the retainined proof in the third case to read the syntactic complexity of the set of decided pairs off the first-order form of the forcing question. That lemma is counted by a separate card of this lane. The number is the article's own section-wise number for it, counted over the numbered environments of section 6.\par\medskip
\begin{definition}
Let $(f, \sigma, X)$ be a condition and $\varphi(G)$ be a $\Sigma^0_1$-formula.
Let $(f, \sigma, X) \qvdash \varphi(G)$ hold if and only if for every relation $R \in \R_{b_f^{+}}$, there is a finite $g_R$-free set $\rho \subseteq X$ such that $\varphi(\sigma \cup \rho)$ holds.
\end{definition}

We say that a condition $p$ \emph{forces} a formula $\varphi(G)$ if $\varphi(G_\F)$ holds for every sufficiently generic filter~$\F$ containing~$p$.

\begin{lemma}\label[lemma]{lem:fs-diagonalization}
Let $p = (f, \sigma, X)$ be a condition and $\Phi_e$ be a Turing functional.
There is an extension of~$p$ forcing $\Phi_e^G \neq D$.
\end{lemma}

\begin{proof}
Let $U = \{ (x, v) \in \NN \times 2 : p \qvdash \Phi^G_e(x) \downarrow = v \}$. We have three cases:
\begin{itemize}
    \item Case 1: $(x, 1-D(x)) \in U$ for some~$x \in \NN$. By \Cref{lem:fs-forcing-question-spec}, there is an extension of~$p$ forcing $\Phi^G_e(x)\downarrow = 1-D(x)$, hence forcing $\Phi_e^G \neq D$.
    \item Case 2: $(x, D(x)) \not \in U$ for some~$x \in \NN$. By \Cref{lem:fs-forcing-question-spec}, there is an extension of~$p$ forcing $\neg (\Phi^G_e(x)\downarrow = D(x))$, hence forcing $\Phi_e^G \neq D$.
    \item Case 3: $U$ is the graph of the characteristic function of~$D$. By \Cref{lem:fs-forcing-question-sigma01}, the set $U$ is $\Sigma^0_1(X)$, so $D \leq_T X$, contradiction.
\end{itemize}
\end{proof}

\end{document}
